\begin{table}[H]
\centering
\scriptsize
\caption{Taxonomía de FungiTastic y número de especies presentes en el dataset.}
\label{tab:taxmap}
\begin{tabularx}{\linewidth}{l Y Y Y Y r}
\toprule
\textbf{Género} & \textbf{Familia} & \textbf{Orden} & \textbf{Clase} & \textbf{Filo} & \textbf{Especies}\\
\midrule
\textit{Amanita} & Amanitaceae & Agaricales & Agaricomycetes & Basidiomycota & 19\\
\textit{Auricularia} & Auriculariaceae & Auriculariales & Agaricomycetes & Basidiomycota & 2\\
\textit{Ganoderma} & Polyporaceae & Polyporales & Agaricomycetes & Basidiomycota & 5\\
\textit{Geastrum} & Geastraceae & Geastrales & Agaricomycetes & Basidiomycota & 9\\
\textit{Hericium} & Hericiaceae & Russulales & Agaricomycetes & Basidiomycota & 3\\
\textit{Lycoperdon} & Lycoperdaceae & Agaricales & Agaricomycetes & Basidiomycota & 10\\
\textit{Morchella} & Morchellaceae & Pezizales & Pezizomycetes & Ascomycota & 3\\
\textit{Pleurotus} & Pleurotaceae & Agaricales & Agaricomycetes & Basidiomycota & 3\\
\textit{Trametes} & Polyporaceae & Polyporales & Agaricomycetes & Basidiomycota & 7\\
\textit{Xylaria} & Xylariaceae & Xylariales & Sordariomycetes & Ascomycota & 5\\
\bottomrule
\end{tabularx}
\source{data/taxonomy/v002/genus_taxonomy.csv}
\end{table}
